% ===========================================================================
%  Every \usepackage of the article documents lives here.
%
%  Loaded first by preamble.tex. style.tex and commands.tex assume everything
%  below is already available -- this file therefore has to stay self-contained
%  and free of layout decisions.
% ===========================================================================

% ---------------------------------------------------------------------------
% Build switches -- these come first because every later file keys off them.
%
%   \npstyle = screen  (default)  coloured links and info boxes, screen header
%            = print              grayscale, two-sided, print-optimised
%            = plain              no fancy header/footer at all (short handouts)
%
% A document selects one by defining \npstyle before loading the preamble.
% ---------------------------------------------------------------------------

\usepackage{etoolbox}               % \ifdefstring, \patchcmd
\usepackage{ifthen}                 % \ifthenelse, used by the macros in commands.tex

\providecommand{\npstyle}{screen}
\newif\ifnpprint \ifdefstring{\npstyle}{print}{\npprinttrue}{\npprintfalse}
\newif\ifnpfancy \ifdefstring{\npstyle}{plain}{\npfancyfalse}{\npfancytrue}

% ---------------------------------------------------------------------------
% Who wrote this. Loaded here because this is the one file every document
% pulls in, so \npauthor / \npauthorurl / \npauthormail are always available.
% ---------------------------------------------------------------------------
% ===========================================================================
%  Who wrote this and where it lives.
%
%  These three macros are the ONLY place the author's identity appears in the
%  layout (running heads, footers, title pages, slide footers).  If you fork
%  this repository, change them here and every document follows.
% ===========================================================================

\providecommand{\npauthor}{Hendrik Möller}
\providecommand{\npauthorurl}{https://github.com/Hendrik-code/numprog_script}
\providecommand{\npauthormail}{hendrik.moeller[AT]tum.de}


% ---------------------------------------------------------------------------
% Encoding, fonts and language
% ---------------------------------------------------------------------------
\usepackage[T1]{fontenc}            % proper hyphenation of words with umlauts
\usepackage[utf8]{inputenc}         % source files are UTF-8
\usepackage{lmodern}                % scalable Computer-Modern-like font
\usepackage{mathpazo}               % Palatino for the title pages
\usepackage[ngerman]{babel}         % German hyphenation and captions
\usepackage[autostyle, german=guillemets]{csquotes}
\MakeOuterQuote{"}                  % lets you type "..." for German quotes
\usepackage{relsize}

% ---------------------------------------------------------------------------
% Page layout
% ---------------------------------------------------------------------------
\usepackage{fancyhdr}               % configured in style.tex
\usepackage{multicol}
\usepackage{float}                  % [H] float placement
\usepackage{placeins}               % \FloatBarrier
\usepackage{framed}                 % the coloured info boxes
\usepackage{mdframed}
\usepackage{rotating}

% ---------------------------------------------------------------------------
% Maths
% ---------------------------------------------------------------------------
\usepackage{amsmath, accents}
\usepackage{amssymb}
\usepackage{cancel}
\usepackage{centernot}              % strike through long symbols such as ==>
\usepackage{esdiff}                 % derivative operator
\usepackage{nicefrac}
\usepackage{fp}                     % fixed point arithmetic in macros

% ---------------------------------------------------------------------------
% Tables, colour and symbols
% ---------------------------------------------------------------------------
\usepackage{booktabs}
\usepackage{colortbl}
\usepackage{xcolor}
\usepackage{enumerate}
\usepackage{pifont}                 % \ding checkmarks and crosses
\usepackage{tikzsymbols}            % \Smiley, \Cooley, \Laughey, ...
\tikzsymbolsset{global-scale=1.2}

% ---------------------------------------------------------------------------
% Graphics and plots
%
% pgfplots pulls in tikz; it MUST be loaded here because style.tex
% calls \pgfplotsset.  (Documents used to load only the style file and failed
% with "Undefined control sequence" on \pgfplotsset -- that is why there is a
% single entry point now.)
% ---------------------------------------------------------------------------
\usepackage{graphicx}
\usepackage{tikz}
\usetikzlibrary{datavisualization}
\usetikzlibrary{datavisualization.formats.functions}
\usepackage{pgfplots}
\pgfplotsset{compat=1.16}
\usetikzlibrary{pgfplots.colorbrewer}

% ---------------------------------------------------------------------------
% Code listings
% ---------------------------------------------------------------------------
\usepackage{listings}
\lstset{
    frame            = tb,
    language         = Python,
    aboveskip        = 3mm,
    belowskip        = 3mm,
    showstringspaces = false,
    columns          = flexible,
    basicstyle       = {\small\ttfamily},
    numbers          = none,
    breaklines       = true,
    breakatwhitespace= true,
    tabsize          = 3,
}

% ---------------------------------------------------------------------------
% Programming helpers used by commands.tex
% ---------------------------------------------------------------------------
\usepackage{environ}                % \NewEnviron
\usepackage{catchfilebetweentags}   % pulling snippets into the complete script
\usepackage{nccfoots}               % \Footnotetext
\usepackage{verbatim}
\usepackage{hyperref}               % keep last: it patches other packages
